\scanpage{210}
\mspage{-2.10-}
% archive sheet number: 209
(i) $\handwritten{\psi_{\unsure{x|\chi}^*}}$ is not proper if and only if for some non-trivial 1-PS $\handwritten{\lambda}$ of $G$, $\handwritten{\sigma}^*(\handwritten{\lambda(\alpha)},x^*)$ has a specialization when $\handwritten{\alpha\longrightarrow}0$.

(ii) $(0)$ is in the closure of $\unsure{O|0}(x^*)$ if and only if for some 1-PS $\handwritten{\lambda}$ of $G$, $\handwritten{\sigma}^*(\handwritten{\lambda(\alpha)},x^*)$ specializes to $(0)$ when $\handwritten{\alpha\longrightarrow}0$.

\uline{Proof}: The two ``if'' statements are obvious. Next consider the ``only if'' in (i): assume $\handwritten{\psi_{x^*}}$ is not proper. Then by the valuative criterion for properness (EGA 2, \S7), \unsure{there-is|there is} a $K$-valued point $\handwritten{\phi}:\operatorname{Spec}(K)\handwritten{\longrightarrow}G$ such that $\handwritten{\phi}$ is not induced by an $R$-valued point of $G$, but such that $\handwritten{\psi_{x^*}\circ\phi}$ is an $R$-valued point of $\handwritten{\mathbb{A}}^n$. By Iwahori's theorem, then, we know that $\handwritten{\phi}$ has the form:
\[
\handwritten{\phi}=\handwritten{\psi}_1\cdot\handwritten{\langle\lambda\rangle}\cdot\handwritten{\psi}_2
\]
for some 1-PS $\handwritten{\lambda}$ of $G$, and for $\handwritten{\psi}_1,\handwritten{\psi}_2\handwritten{\unsure{\in|\epsilon}}G(R)$. Moreover, $\handwritten{\lambda}$ is non-trivial since $\handwritten{\phi}$ is not itself an $R$-valued point of $G$. Let $b_i\handwritten{\in}G_k$ be $k$-valued point $\handwritten{\unsure{\varpi|\overline{\omega}}}(\handwritten{\psi}_i)$ obtained from $\handwritten{\psi}_i$ when $t$ specializes to $0$. By choosing suitable coordinates in $\handwritten{\mathbb{A}}^n$, we may assume that the action of $b_2^{-1}\cdot\handwritten{\lambda}\cdot b_2$ on $\handwritten{\mathbb{A}}^n$ is diagonalized. Suppose $b_2^{-1}\cdot\handwritten{\lambda(\alpha)}\cdot b_2$ acts via the matrix $(\handwritten{\alpha}^{\handwritten{\unsure{r|n}_i}}\cdot\handwritten{\delta}_{ij})$, for $\handwritten{\alpha\in}(\handwritten{\mathbb{G}}_m)_k$. Now
\[
\begin{aligned}
\handwritten{\psi_{x^*}\circ\phi}&=\handwritten{\sigma}^*(\handwritten{\phi},x^*)\\
&=\handwritten{\sigma}^*\handwritten{\bigl(}(\handwritten{\psi}_1\cdot b_2)\cdot(b_2^{-1}\cdot\handwritten{\langle\lambda\rangle}\cdot b_2)\cdot(b_2^{-1}\cdot\handwritten{\psi}_2),x^*\handwritten{\bigr)}
\end{aligned}
\]
(Here $k$-valued points of $G$ and of $\handwritten{\mathbb{A}}^n$ are identified with the $K$-valued points obtained by the base extension $K\handwritten{\supset}k$).
